\chapter{未確認事項と完成に向けた確認順序}\label{app:gaps}
本付録は執筆用の不足一覧であり，提出稿では解消した項目を反映して整理する．赤字の要確認事項と，性能上すでに観測された未達を分ける．例えばONB段階の検知不足は，数値をまだ集計していない空欄ではなく，現在の結果が示した到達限界である．その未達を消すために目的を全域RMSEだけへ変更しない．

\section{出典・記述の完成に必要な事項}
\begin{longtable}{L{13mm}L{46mm}L{70mm}}
\caption{本文の要確認IDと必要な根拠．}\\
\toprule ID & 不足 & 完了に必要な確認 \\ \midrule\endfirsthead
\toprule ID & 不足 & 完了に必要な確認 \\ \midrule\endhead
D01 & データ取得者・提供元 & 原実験ノートとWAVを対応づけ，実施者・提供元を確認する \\
D02 & 著者等の担当範囲 & 計測，ラベル算出，前処理，実装，実行，解析・考察の担当を分ける \\
D03 & 対象日の実験・収録条件 & 寸法，水温・圧力，センサ位置，ゲイン，校正を原資料で確認する \\
D04 & 熱流束ラベルの算出と誤差 & CSVを原電気計測値へ遡り，面積，平均時間，不確かさを確認する \\
D05 & ONB判定の原根拠 & 使用値の確認記録に加え，原観察，判定基準，分解能を示す \\
D06 & 学習時の環境版 & 主実行に対応するソフトウェア版・ハードウェア記録を確定する \\
L01 & 先行研究本文の確認 & 対象・分割・雑音・ONB評価の比較表を作り，未読部分を補う \\
L02 & 物理考察の照合 & 手塚最終修論と原実験条件を対応させ，音源候補の根拠を整理する \\
L03 & IGを追加する場合の根拠 & 対象モデル・入力・基準・収束を揃える．現在の中心結論の必須図ではない \\
F01 & 題目・提出書式 & 所定要領に表紙，余白，肩書，年月，引用形式を照合する \\
F02 & 要旨・用語の最終校閲 & 指導後に和英要旨と用語，本文・図表との整合を確認する \\
F03 & 謝辞 & 氏名・協力内容を本人が確認して記入する \\
\bottomrule
\end{longtable}
\need{F02}{和英要旨は現在の限定した結論に合わせて作成した．正式題目と指導後の本文に合わせ，専門用語，英文および提出時の字数条件を最終確認する．}

\section{強い結論を採る場合にのみ必要な追加証拠}
\begin{longtable}{L{13mm}L{45mm}L{71mm}}
\caption{追加検証の必要性を主張の強さと対応づける．}\\
\toprule ID & 採りたい強い主張 & 必要となる識別根拠 \\ \midrule\endfirsthead
\toprule ID & 採りたい強い主張 & 必要となる識別根拠 \\ \midrule\endhead
V01 & 200 epochsが最良・学習十分 & 学習側のvalidation曲線と事前に定めたepoch選択 \\
V02 & 未知録音・未知日でも優位 & 開発に未参照の独立データ，分割単位，固定した選択規則 \\
V03 & 入力依存差が雑音利得の原因・帯域の物理同定 & 表現・学習器を統制した比較，原実験条件と独立観測の照合 \\
V04 & ONB発生段階の十分な早期検知 & 許容FN/FPと対象母集団を定めた評価．秒単位を言うなら発生時刻の正解も必要 \\
\bottomrule
\end{longtable}
\need{V04}{現在の30区間の結果は開始段階の検知未達を示す．到達熱流束を実時間遅延や実機の安全余裕へ換算できる根拠はない．現在の限定した結論では未達として記述する．}

\section{本文を完成させるための順序}
まずD01--D05により，データと正解の根拠および研究上の立場を明確にする．次にL01--L02を補い，目的と物理考察の文章を根拠へ対応づける．併せて，第1章の目的と第6章の結論を読み合わせ，全域回帰，雑音耐性，早期検知，説明性という語の強さが揃っていることを確認する．

その上で，指導により必要とされる達成範囲を定め，現在の限定した結論を維持するか，V01--V04のどの主張を強めるかを判断する．追加実行を行う場合は，比較対象，必要な独立性，改善対象，終了条件を先に記録する．モデル探索を広げることや過去の完了確認を繰り返すことを，執筆の前提にはしない．

最後に，本文で使う図表・引用の対応，実験資料の出典，謝辞および提出様式を整える．本稿は全章の文章化を行った状態であり，データ出典等の未確認，ONB検知上の未達，指導・提出要領による修正が残る．修論の合格基準や教授の了承を満たした完成稿という意味ではない．
